%! Dimensions of the Four Subspaces — Unit 3, Lesson 3.5 — Guided Notes
\documentclass[10pt]{article}
\usepackage{linalg-article}
\usepackage{linalg-boxes}
\newcommand{\TallMath}[1]{$\displaystyle #1\rule[-1.4em]{0pt}{3.2em}$}
\newcommand{\vv}{\mathbf{v}}
\newcommand{\ww}{\mathbf{w}}
\newcommand{\xx}{\mathbf{x}}
\newcommand{\yy}{\mathbf{y}}
\newcommand{\cc}{\mathbf{c}}
\newcommand{\bb}{\mathbf{b}}
\newcommand{\svec}{\mathbf{s}}
\newcommand{\zero}{\mathbf{0}}
\newcommand{\T}{\mathsf{T}}

\begin{document}

\pageheader{Unit 3, Lesson 3.5}{Guided Notes}
\namedateperiod

\vspace{0.06in}

\begin{objectivebox}
\textbf{By the end of this lesson, I will be able to\dots}
\begin{itemize}\setlength{\itemsep}{1pt}
  \item name the \emph{four fundamental subspaces} of $A$ and say which room ($\mathbb{R}^m$ or $\mathbb{R}^n$) each lives in;
  \item give every dimension from the single number $r$: $\dim C(A)=\dim C(A^{\T})=r$, $\dim N(A)=n-r$, $\dim N(A^{\T})=m-r$;
  \item read a \emph{basis} for all four off one reduction, and check $r+(n-r)=n$ and $r+(m-r)=m$.
\end{itemize}
\end{objectivebox}

\vspace{0.04in}

\begin{vocabbox}
\small
Fill in each term as we build it together.
\vspace{2pt}
\termblanklong{Row space $C(A^{\T})$}
\termblanklong{Left nullspace $N(A^{\T})$}
\termblanklong{Row rank $=$ column rank}
\termblanklong{Counting rules}
\end{vocabbox}

\begin{hookbox}
The robot arm connects \textbf{two rooms}: a lever setting is an input $\xx$ in $\mathbb{R}^n$, and the tip
lands at the output $A\xx$ in $\mathbb{R}^m$.
\begin{itemize}\setlength{\itemsep}{1pt}
  \item Which input settings leave the tip \emph{fixed} (do nothing)? That subspace is called the \writeline
  \item Which output positions can the tip actually \emph{reach}? That subspace is called the \writeline
\end{itemize}
\end{hookbox}

\begin{notesbox}{1. Two Rooms, Four Subspaces}
An $m\times n$ matrix takes inputs $\xx\in\mathbb{R}^n$ to outputs $A\xx\in\mathbb{R}^m$ --- so it lives
between an \textbf{input room} $\mathbb{R}^n$ and an \textbf{output room} $\mathbb{R}^m$. Each room splits
into \emph{two} fundamental subspaces:
\begin{center}\small
\renewcommand{\arraystretch}{1.3}
\begin{tabular}{@{}l l l@{}}
\textbf{Subspace} & \textbf{What it holds} & \textbf{Room} \\
\hline
Column space $C(A)$ & all combinations of the columns (reachable outputs) & $\mathbb{R}^{\blank{0.7cm}}$ \\
Left nullspace $N(A^{\T})$ & all $\yy$ with $A^{\T}\yy=\zero$ (rows combine to $\zero$) & $\mathbb{R}^{\blank{0.7cm}}$ \\
Row space $C(A^{\T})$ & all combinations of the rows & $\mathbb{R}^{\blank{0.7cm}}$ \\
Nullspace $N(A)$ & all $\xx$ with $A\xx=\zero$ (do-nothing inputs) & $\mathbb{R}^{\blank{0.7cm}}$ \\
\end{tabular}
\end{center}
The column space and left nullspace live in the \textbf{output} room $\mathbb{R}^m$; the row space and
nullspace live in the \textbf{input} room $\mathbb{R}^n$. The row space is just ``the column space of
$A^{\T}$,'' and the left nullspace is ``the nullspace of $A^{\T}$'' --- transpose, and you are back to tools
from Lessons~1.3 and~3.2.
\end{notesbox}

\begin{notesbox}{2. Every Dimension Comes from $r$}
The rank $r$ (the number of pivots) fixes all four dimensions:
\[
  \dim C(A)=\blank{0.7cm}, \qquad \dim C(A^{\T})=\blank{0.7cm}, \qquad
  \dim N(A)=\blank{1.4cm}, \qquad \dim N(A^{\T})=\blank{1.4cm}.
\]
The headline is \textbf{row rank $=$ column rank}: the number of independent rows equals the number of
independent columns --- \emph{both} are $r$. So $\dim C(A)=\dim C(A^{\T})=r$ even when the two spaces live
in different-sized rooms. \textbf{Two counting rules} follow, one per room:
\[
  \underbrace{r+(n-r)}_{\text{input room}}=n, \qquad\qquad \underbrace{r+(m-r)}_{\text{output room}}=m.
\]
Each room is split into ``directions that do something'' ($r$-dimensional) plus ``directions that vanish.''
\end{notesbox}

\begin{notesbox}{3. All Four Bases from One Reduction}
Reduce $A$ \emph{once} and read off everything. Take
$A=\begin{bmatrix} 1 & 1 & 2 \\ 2 & 1 & 3 \\ 3 & 2 & 5 \end{bmatrix}$ ($m=3$, $n=3$):
\[
  \begin{bmatrix} 1 & 1 & 2 \\ 2 & 1 & 3 \\ 3 & 2 & 5 \end{bmatrix}
  \xrightarrow{\substack{R_2-2R_1\\ R_3-3R_1}}
  \begin{bmatrix} 1 & 1 & 2 \\ 0 & -1 & -1 \\ 0 & -1 & -1 \end{bmatrix}
  \xrightarrow{\substack{R_3-R_2,\ -R_2\\ R_1-R_2}}
  \begin{bmatrix} 1 & 0 & 1 \\ 0 & 1 & 1 \\ 0 & 0 & 0 \end{bmatrix}.
\]
Pivots in columns $1,2$, so $r=\blank{0.7cm}$. Now read the four bases:
\begin{itemize}\setlength{\itemsep}{2pt}
  \item \textbf{$C(A)$} (in $\mathbb{R}^3$): the \emph{pivot columns of $A$} $\begin{bsmallmatrix} 1\\2\\3 \end{bsmallmatrix},\begin{bsmallmatrix} 1\\1\\2 \end{bsmallmatrix}$; $\dim=r=\blank{0.6cm}$.
  \item \textbf{$C(A^{\T})$} (in $\mathbb{R}^3$): the \emph{nonzero rows} of the reduced form $(1,0,1),(0,1,1)$; $\dim=r=\blank{0.6cm}$.
  \item \textbf{$N(A)$} (in $\mathbb{R}^3$): the special solution $\svec=\begin{bsmallmatrix}\rule{0pt}{1.0em}\blank{0.55cm}\\ \blank{0.55cm}\\ \blank{0.55cm}\end{bsmallmatrix}$ (column~3 free); $\dim=n-r=\blank{0.6cm}$.
  \item \textbf{$N(A^{\T})$} (in $\mathbb{R}^3$): the row-combination giving the zero row. Here row${}_1+{}$row${}_2-{}$row${}_3=\zero$, so $\yy=\begin{bsmallmatrix} 1\\1\\-1 \end{bsmallmatrix}$; $\dim=m-r=\blank{0.6cm}$.
\end{itemize}
Checks: $r+(n-r)=\blank{0.6cm}=n$ and $r+(m-r)=\blank{0.6cm}=m$.
\end{notesbox}

\begin{notesbox}{4. The Two Nullspaces Name the Two Redundancies}
The special solution $\svec=\begin{bsmallmatrix} -1\\-1\\1 \end{bsmallmatrix}$ says $A\svec=\zero$, i.e.\
$-\mathbf{a}_1-\mathbf{a}_2+\mathbf{a}_3=\zero$: \textbf{column}~3 is redundant
($\mathbf{a}_3=\mathbf{a}_1+\mathbf{a}_2$). The left-nullspace vector
$\yy=\begin{bsmallmatrix} 1\\1\\-1 \end{bsmallmatrix}$ says $A^{\T}\yy=\zero$, i.e.\
$\text{row}_1+\text{row}_2-\text{row}_3=\zero$: \textbf{row}~3 is redundant
(row${}_3=$ row${}_1+$ row${}_2$). Same idea, two directions:
\begin{center}\small
\begin{tabular}{@{}l l@{}}
Nullspace vector $\to$ a \blank{2.2cm} dependency \quad&\quad Left-nullspace vector $\to$ a \blank{2.2cm} dependency
\end{tabular}
\end{center}
\textbf{Next (3.6):} we assemble these four subspaces --- and their dimensions $r,r,n-r,m-r$ --- into the
``big picture'' of a matrix (the Fundamental Theorem of Linear Algebra), and see how the two pairs sit at
right angles.
\end{notesbox}

\begin{practicebox}
Show your reasoning. Use $r$, $m$, and $n$.
\begin{enumerate}[leftmargin=1.6em, itemsep=2pt]
  \item A matrix $A$ is $4\times 6$ with rank $r=3$. Give all four dimensions: $\dim C(A)$, $\dim C(A^{\T})$, $\dim N(A)$, $\dim N(A^{\T})$. \writeline
  \item For $A=\begin{bmatrix} 1 & 2 & 3 \\ 2 & 4 & 6 \end{bmatrix}$, find $r$, then give a basis for the row space and for the left nullspace, with their dimensions. \writeline
  \item In one sentence: what does a left-nullspace vector tell you about the \emph{rows} of $A$? \writeline
\end{enumerate}
\end{practicebox}

\end{document}
